%%=============================================================================
%% 第 2 章 相关工作
%%=============================================================================

\chapter{相关工作}
\label{chap:related_work}

% 引言
运动想象脑机接口（Motor Imagery Brain-Computer Interface, MI-BCI）作为非侵入式脑机接口的重要分支，近年来在康复医疗、人机交互等领域展现出广阔的应用前景\cite{pfurtscheller1999event}。然而，EEG 信号的非平稳性和个体差异性使得 MI-BCI 系统的性能仍面临诸多挑战。其中，时间窗选择作为影响特征提取质量和分类性能的关键因素，已成为该领域的研究热点之一\cite{lotte2011regularizing, yang2015time}。与此同时，强化学习作为一种自适应决策框架，在 BCI 系统的优化中展现出巨大潜力\cite{zhang2021mars}。本章将系统回顾运动想象 BCI 系统的基本原理、EEG 特征提取与分类方法、时间窗优化技术以及强化学习在 BCI 中的应用现状，为本研究的方法设计奠定理论基础。

本章内容组织如下：第\ref{sec:mi_bci_system}节介绍运动想象脑机接口系统的基本架构与神经生理基础；第\ref{sec:feature_classification}节阐述 EEG 特征提取与分类的核心方法，重点讨论 CSP-SVM 经典流水线；第\ref{sec:time_window_optimization}节系统梳理时间窗优化方法，包括固定时间窗策略与各类自适应优化技术；第\ref{sec:rl_in_bci}节回顾强化学习基础理论及其在 BCI 领域的应用进展；最后，基于文献综述提出本研究的核心研究问题。

\section{运动想象脑机接口系统}
\label{sec:mi_bci_system}

\subsection{运动想象的神经生理基础}
\label{subsec:neuro_basis}

运动想象（Motor Imagery, MI）是指个体在头脑中模拟执行特定动作的心理过程\cite{pfurtscheller1999event}。当用户想象肢体运动时，大脑感觉运动皮层会产生事件相关去同步化（Event-Related Desynchronization, ERD）现象，即$\mu$频段（8--13 Hz）和$\beta$频段（13--30 Hz）的脑电功率显著降低\cite{pfurtscheller1999event}。ERD/ERS 现象的空间分布具有任务特异性，为基于运动想象的 BCI 分类提供了神经生理学基础\cite{pfurtscheller1999event, mcfarland1997spatial}。然而，ERD 的时间动态特性存在显著的个体间差异，使得固定时间窗策略难以适应所有被试\cite{zhang2019temporally, lotte2011regularizing}。

\subsection{BCI 系统的基本架构}
\label{subsec:bci_architecture}

典型的运动想象 BCI 系统包含数据采集、预处理、特征提取、特征选择、分类和后处理等模块\cite{tangermann2012review, lotte2007review}。其中，时间窗选择属于特征选择子任务，直接影响特征提取的质量和分类性能。

\section{EEG 特征提取与分类方法}
\label{sec:feature_classification}

在运动想象 BCI 系统中，特征提取和分类是核心模块\cite{lotte2007review}。本研究采用经典的 CSP-SVM 流水线作为基础分类器，以下介绍相关方法。

\textbf{共空间模式（CSP）}。CSP 是运动想象 BCI 中最常用的特征提取方法\cite{blankertz2008optimizing, mcfarland1997spatial}。该方法通过寻找最优空间滤波器，使得两类任务的方差差异最大化。

给定两类 EEG 数据，CSP 首先计算各类别的平均协方差矩阵：
\begin{equation}
\mathbf{C}_c = \frac{1}{N_c} \sum_{i \in \text{class}_c} \mathbf{X}_i \mathbf{X}_i^\top, \quad c = 1, 2
\end{equation}
其中，$\mathbf{X}_i \in \mathbb{R}^{C \times T}$ 为第 $i$ 个试次的 EEG 数据，$C$ 为通道数，$T$ 为时间点数，$N_c$ 为类别 $c$ 的试次数\cite{blankertz2008optimizing}。

然后，CSP 求解广义特征值问题：
\begin{equation}
\mathbf{C}_1 \mathbf{w} = \lambda (\mathbf{C}_1 + \mathbf{C}_2) \mathbf{w}
\end{equation}
得到特征值 $\lambda$ 和对应的特征向量 $\mathbf{w}$\cite{blankertz2008optimizing}。选择最大和最小的 $m$ 个特征值对应的特征向量构建空间滤波器 $\mathbf{W} \in \mathbb{R}^{2m \times C}$。对每个试次，提取对数方差特征：
\begin{equation}
\mathbf{f} = \log \left( \text{var}(\mathbf{W} \mathbf{X}) \right)
\end{equation}
CSP 计算高效，在二分类任务中表现优异，但对噪声敏感且仅适用于二分类场景\cite{lotte2011regularizing}。

\textbf{支持向量机（SVM）}。SVM 通过寻找最大间隔超平面实现分类，在小样本场景下表现稳健\cite{lotte2007review}。线性 SVM 的目标函数为：
\begin{equation}
\min_{\mathbf{w}, b} \frac{1}{2} \|\mathbf{w}\|^2 + C \sum_{i=1}^N \xi_i
\end{equation}
其中，$\mathbf{w}$ 为超平面法向量，$b$ 为偏置，$C$ 为正则化参数，$\xi_i$ 为松弛变量。SVM 可通过正则化参数控制过拟合风险，适合 BCI 数据稀缺场景。尽管如此，其性能的上限仍然高度依赖于输入特征的质量，而这正是由时间窗的选择所决定的。

除 CSP-SVM 外，研究者还提出了 PSD、小波变换等特征提取方法，以及 LDA、深度学习等分类算法\cite{lotte2007review}。然而，这些方法的性能均受时间窗选择的影响，这为本研究的时间窗优化方法提供了应用空间。

\section{时间窗优化方法}
\label{sec:time_window_optimization}

\subsection{固定时间窗方法}
\label{subsec:fixed_window}

早期 BCI 研究通常采用固定时间窗策略，即根据先验知识选择固定的时间段用于特征提取\cite{lotte2007review}。常见做法包括：

（1）\textbf{刺激后固定延迟}。选择刺激呈现后固定延迟的时间窗，如 0.5--3.5 秒\cite{tangermann2012review}。该方法假设所有被试的 ERD 在相同时间范围内出现，但忽略了个体差异性。

（2）\textbf{基于 ERD 峰值}。通过分析被试的平均 ERD 时间曲线，选择 ERD 最强的时间段\cite{pfurtscheller1999event}。该方法需要额外的校准实验，且 ERD 峰值可能随任务状态变化。

（3）\textbf{经验时间窗}。根据文献报道或研究者经验选择时间窗，如 2--4 秒\cite{lotte2007review}。该方法实现简单，但缺乏理论依据。

固定时间窗方法的优势在于实现简单、计算高效。然而，其局限性在于：（1）无法适应个体间差异，不同被试的最优时间窗可能不同；（2）无法适应个体内差异，同一被试在不同试次或不同 session 的 ERD 时间动态可能变化；（3）需要手动调参，依赖研究者经验\cite{lotte2011regularizing}。

\subsection{基于优化的时间窗选择}
\label{subsec:optimization_based}

为克服固定时间窗的局限性，研究者提出了基于优化的自适应时间窗选择方法。

\subsubsection{贝叶斯优化方法}

Yang 等人（2015）提出使用贝叶斯优化（Bayesian Optimization）搜索最优时间窗参数\cite{yang2015time}。该方法构建高斯过程代理模型，近似时间窗参数与分类准确率的映射关系，并通过采集函数（如期望改进 EI）平衡探索与利用。实验结果表明，贝叶斯优化在 BCI Competition IV 2a 数据集上显著优于固定时间窗方法\cite{yang2015time}。然而，贝叶斯优化的计算复杂度随搜索次数增加而增长（高斯过程推断复杂度为 $O(n^3)$），且难以处理高维参数空间（如同时优化时间窗、频段和 CSP 分量数）。这种计算负担对于需要快速校准的在线 BCI 系统而言，是难以接受的。

\subsubsection{网格搜索与随机搜索}

最直接的方法是遍历所有可能的时间窗组合（网格搜索），或随机采样时间窗参数（随机搜索），选择验证集准确率最高的参数。该方法实现简单，但计算效率低，尤其在参数空间较大时。

\subsubsection{遗传算法}

遗传算法（Genetic Algorithm, GA）通过模拟自然选择过程搜索最优时间窗参数。该方法可处理离散和连续参数，但收敛速度慢，且需要设计合适的编码和变异策略。

\subsubsection{稀疏组空间模式}

Zhang 等人（2019）提出时间约束稀疏组空间模式（Temporally Constrained Sparse Group Spatial Patterns）方法，通过正则化项鼓励时间窗的稀疏性和连续性\cite{zhang2019temporally}。该方法将时间窗选择嵌入 CSP 优化过程，实现联合优化。然而，该方法增加了 CSP 的计算复杂度，且正则化参数需要手动调优。

\subsection{时间窗优化的评估标准}
\label{subsec:evaluation_metrics}

时间窗优化方法的评估通常采用以下指标\cite{tangermann2012review, lotte2007review}：

（1）\textbf{分类准确率}。使用最优时间窗提取特征后，分类器在测试集上的准确率。这是最直接的评估指标。

（2）\textbf{Kappa 系数}。Cohen's Kappa 系数考虑了随机猜测的影响，定义为：
\begin{equation}
\kappa = \frac{p_o - p_e}{1 - p_e}
\end{equation}
其中，$p_o$ 为观测准确率，$p_e$ 为随机猜测期望准确率\cite{cohen1960coefficient}。Kappa 系数在多分类场景下更具可比性。

（3）\textbf{计算效率}。包括搜索时间（找到最优时间窗所需的时间）和推理延迟（使用最优时间窗进行分类的时间）\cite{yang2015time}。

（4）\textbf{稳定性}。跨被试或跨 session 的标准差，反映方法对不同条件的鲁棒性\cite{lotte2011regularizing}。

本研究采用上述指标全面评估所提方法的性能，并与贝叶斯优化和 DQN 方法进行对比。

\section{强化学习在 BCI 中的应用}
\label{sec:rl_in_bci}

\subsection{强化学习基础}
\label{subsec:rl_basics}

强化学习（Reinforcement Learning, RL）是机器学习的一个分支，关注智能体如何通过与环境交互学习最优策略\cite{sutton2018reinforcement}。强化学习问题通常建模为马尔可夫决策过程（MDP），形式化定义为五元组 $(\mathcal{S}, \mathcal{A}, \mathcal{R}, \mathcal{P}, \gamma)$\cite{sutton2018reinforcement}。

\textbf{Q-Learning} 是一种经典的无模型强化学习算法，由 Watkins 和 Dayan 于 1992 年提出\cite{watkins1992q}。Q-Learning 通过维护状态 - 动作价值表（Q 表），使用 TD 误差迭代更新 Q 值，在有限状态空间下保证收敛到最优策略\cite{watkins1992q}。由于其无需环境模型、实现简单且在低维状态空间下高效，Q-Learning 在资源受限场景中具有独特优势。

\subsection{深度 Q 网络}
\label{subsec:deep_q_network}

深度 Q 网络（Deep Q-Network, DQN）是强化学习与深度学习结合的典型代表。Mnih 等人（2015）提出 DQN，使用深度神经网络近似 Q 函数\cite{mnih2015human}。DQN 通过最小化时序差分损失函数训练网络，并引入经验回放（Experience Replay）和目标网络（Target Network）两项关键技术提升训练稳定性\cite{mnih2015human}。DQN 在 Atari 游戏等复杂任务中取得了突破性进展，但其模型复杂、需要大量训练数据且可解释性较差。

后续研究提出了多种 DQN 改进方法，包括 Double Q-Learning 解决 Q 值过高估计问题\cite{van2016deep, hasselt2010double}、Dueling 网络架构将 Q 函数分解为状态价值函数和优势函数\cite{wang2016dueling}，以及集成多种改进的 Rainbow 算法\cite{hessel2018rainbow}。然而，这些改进进一步增加了模型的复杂度和计算开销。

\subsection{强化学习在 BCI 中的应用}
\label{subsec:rl_bci_applications}

近年来，强化学习被应用于 BCI 系统的多个方面\cite{zhang2021mars}：

\textbf{动态窗口选择}。Chen 等人（2019）提出基于强化学习的动态窗口方法，用于 SSVEP（Steady-State Visually Evoked Potential）识别\cite{chen2019dynamic}。该方法将窗口长度选择建模为 MDP，使用 Q-Learning 学习最优策略，在信息传输率和准确率上均优于固定窗口方法。该研究为时间窗优化提供了方法迁移思路，但针对的是 SSVEP 范式，而非运动想象。

\textbf{多智能体特征选择}。Zhang 等人（2021）提出 MARS（Multiagent Reinforcement Learning for Spatial-spectral and Temporal feature selection）方法，使用多智能体强化学习联合优化 EEG 的空间、频谱和时间特征\cite{zhang2021mars}。MARS 在多个 BCI 数据集上取得了先进性能，但计算复杂度高，需要大量训练数据，且可解释性较差。

\textbf{自适应 BCI 控制}。强化学习还被用于 BCI 系统的自适应控制，如调整分类器参数、选择最优 BCI 范式和优化神经反馈策略\cite{jayaram2016transfer}。

\textbf{迁移学习}。Jayaram 等人（2016）综述了迁移学习在 BCI 中的应用，包括跨被试迁移、跨 session 迁移和跨任务迁移\cite{jayaram2016transfer}。迁移学习可减少新被试的校准时间，但面临领域偏移和负迁移等挑战。

\section{本章小结}
本章系统回顾了运动想象 BCI 系统与时间窗优化方法的相关研究工作。首先，介绍了运动想象的神经生理基础（ERD/ERS 现象）与 BCI 系统的基本架构，阐明了时间窗选择对分类性能的重要影响。其次，阐述了 EEG 特征提取与分类的核心方法，重点讨论了 CSP-SVM 经典流水线及其理论特性。再次，系统梳理了时间窗优化方法的发展历程，从固定时间窗策略到基于贝叶斯优化、遗传算法的自适应方法，分析了各类技术的优势与局限。最后，回顾了强化学习基础理论及其在 BCI 领域的应用进展，特别是 MARS 等多智能体强化学习方法在特征选择中的成功实践。

基于文献综述，本章发现：现有强化学习方法（如 MARS）虽在 BCI 特征选择中表现优异，但其复杂的多智能体设置和深度网络结构带来了巨大的计算开销，难以应用于资源受限的在线 BCI 系统。这引出了本研究的核心研究问题：\textbf{在运动想象时间窗优化这一状态空间有限（136 个状态）的特定任务中，是否可以采用更简洁的表格型 Q-Learning，在避免复杂性的同时，保留强化学习的自适应能力？} 该问题的探索将为轻量级自适应 BCI 系统的设计提供新的思路，也为下一章方法设计奠定理论基础。


